\begin{myexercise}[Arena zur Amerikanisierung der Städte]
Wählen Sie eine folgender Rollen und überlegen Sie sich während dem \href{https://www.srf.ch/play/tv/kulturplatz/video/hochhaeuser-das-leben-in-der-vertikale?urn=urn:srf:video:a98ccb20-f5ad-46ea-8203-2688102543d4}{Dokumentarfilm von SRF}, welche Argumente aus Sicht dieser Personengruppe für oder gegen den Bau weiterer Hochhäuser in der Schweiz sprechen.
\begin{enumerate}
\item StadtplanerIn
\item Immobilien-InvestorIn
\item Soziologe
\item UmweltschützerIn
\item Heimatschutz- und Denkmal-PflegerIn
\item Familien-Mutter / -Vater
\end{enumerate}
Eventuell kommt Ihre Rolle im Film nicht vor. Das macht nichts! Überlegen Sie sich trotzdem, wie Sie sich positionieren würden.

Wir werden basierend auf dem Film eine Debatte zum Thema \enquote{Manhattisierung der Schweiz} durchführen. Dabei führt eine Person die \textbf{Moderation} und eine Person führt \textbf{Protokoll}.

Moderations-Fragen könnten entlang folgender Aspekte aufgebaut sein: 
\begin{itemize}
\item \textbf{Generelles Warm-Up}: braucht es mehr Hochhäuser? Wichtigste Pro- und Kontra-Argumente pro Seite, pro Seite je 2 Personen zu Wort kommen lassen.
\item \textbf{Raumplanung}: Sollten die Menschen einen finanziellen Anreiz erhalten, um in Hochhäusern zu wohnen? Sind Hochhäuser ein sinnvolles Mittel, um Städte zu verdichten?
\item \textbf{Politik}: Sollten Städte die politischen Auflagen für die Planung und den Bau neuer Hochhäuser senken?
\item \textbf{Umwelt-Aspekte}: Sind Hochhäuser energetisch effizienter als herkömmliche Wohngebäude? Führen Sie zu weniger Ressourcenverbrauch?
\item \textbf{Soziales}: Wie verändern Hochhäuser das soziale Zusammenleben in der Stadt?
\item \textbf{Sicherheit}: Wie verändern Hochhäuser das Sicherheitsgefühl der Anwohner? (Schulweg, Nachbarschaft, Anonymität, etc.)
\item \textbf{Ästhetik \& Kultur}: Wie verändern Hochhäuser das Stadtbild?
\end{itemize}

Pro Aspekt sollten maximal 2 Minuten pro Seite eingeplant werden. Es sollte darauf geachtet werden, dass beide Seiten ausgewogen zu Wort kommen.
\end{myexercise}

\begin{mychallenge}[{Ein Lob auf Wolkenkratzer (11.11.2021 NZZ)}]
\begin{figure}[H]
\centering
\colorbox{white}{\includegraphics[trim={1cm 16cm 1cm 1cm},clip,width=.9\linewidth]{"Geografie/Stadtgeografie/Appendix/2021-11-11 Lob auf Wolkenkratzer.pdf"}}
\end{figure}
\end{mychallenge}


\begin{mychallenge}[{Artikel zu geplanten Hochhäusern in Zürich (25.11.2024 NZZ)}]
\begin{figure}[H]
\centering
\colorbox{white}{\includegraphics[trim={1cm 17cm 1cm 1cm},clip,width=.9\linewidth]{"Geografie/Stadtgeografie/Appendix/2024-11-25 Hochhäuser Zürich"}}
\end{figure}
\end{mychallenge}


\begin{mychallenge}[{Essay über Hochhäuser in Zürich (04.09.2023 NZZ Folio)}]
\foreach \pagenum in {1,2} {
    \begin{figure}[H]
    \centering
    \colorbox{white}{\includegraphics[width=.9\linewidth, page=\pagenum]{"Geografie/Stadtgeografie/Appendix/2023-09-04 NZZ Folio Hochhäuser Zürich"}}
    \end{figure}
}
\end{mychallenge}



\begin{mychallenge}[Weitere Unterlagen]
\begin{itemize}
\item \href{https://www.srf.ch/kultur/gesellschaft-religion/verdichtung-in-staedten-wie-kann-zuerich-platz-fuer-zusaetzliche-70-000-menschen-schaffen}{SRF-Beitrag zur Verdichtung}
\item \href{https://www.swissinfo.ch/ger/gesellschaft/serie-raumplanung_warum-in-der-schweiz-kaum-in-die-hoehe-gebaut-wird/45792966}{Artikel: \enquote{Warum in der Schweiz kaum in die Höhe gebaut wird}}
\item \href{https://3d.stzh.ch/appl/3d/zuerich_4d_extern/}{Geplante (Hoch-)Häuser in Zürich}
\end{itemize}
\end{mychallenge}


\newpage
\section{Lernziele}
\begin{todolist}
\item Ich kann den Aufbau \textbf{nordamerikanischer Städte} sowie typische Segregationsformen anhand eines räumlichen Modells beschreiben, indem ich die Begriffe \textit{CBD (Central Business District)}, \textit{Schachbrett-Grundriss} (\textit{Gridiron Pattern}), \textit{Übergangsbereich}, \textit{Auto-Infrastruktur}, \textit{Umland}, \textit{Edge City} sowie \textit{Gated Community} räumlich korrekt einordne.
\item Ich kann die Hauptargumente, welche für und gegen die Verdichtung von Städten sprechen, insbesonderen hinsichtlich dem Bau neuer Wolkenkratzer, aus der Perspektive verschiedener Bevölkerungsgruppen (Familien, Investoren, Stadtplaner, Umweltaktivisten etc.) erläutern.
\end{todolist}

